\documentclass{article}
\usepackage{ngerman}
\usepackage{graphicx}
\usepackage{xfrac}
\usepackage{float}
\usepackage{tabularx}
\usepackage[margin=1in]{geometry}
\usepackage{fancyhdr}
\usepackage{lmodern}

\newcommand{\titel}{Leipzig eSports e.V. Corporate Design Manual}
\newcommand{\autor}{Leipzig eSports e.V.}
\newcommand{\version}{0.2}
\usepackage{xcolor}
\definecolor{LESOrange}{HTML}{FF7929}
\definecolor{LESBlue}{HTML}{2E72C9}
\definecolor{LESYellow}{HTML}{FFD129}

\definecolor{SECLight}{HTML}{b4b8bd}
\definecolor{SECMedium}{HTML}{666666}
\definecolor{SECDark}{HTML}{333333}

\newcommand\cfield[1][black]{\textcolor{#1}{\linebreak\noindent\rule{\textwidth}{1cm}\linebreak}}
\usepackage{ifxetex,fontspec}
\usepackage{titlesec}

\setsansfont{Kanit}
\setmainfont{Titillium Web}

\newfontfamily\headingfont{Kanit}

\titleformat*{\section}{\Large\headingfont}

\setlength{\parindent}{0pt}

\newcommand\alphabet{ABCDEFGHIJKLMNOPQRSTUVWXYZ\\abcdefghijklmnopqrstuvwxyz}

\newcommand\fonttest[1]{\begin{samepage}{\fontsize{18pt}{18pt}\fontspec{#1}\textbf{\alphabet}\\\alphabet}\end{samepage}}


\pagestyle{fancy}
\begin{document}

\begin{titlepage}
\vspace*{\fill}
\title{\titel}
\author{\autor}
\includegraphics[width=\textwidth]{Docs/Logo.eps}
{\let\newpage\relax\maketitle}
\begin{center}
Version \version
\end{center}
\vspace*{\fill}
\end{titlepage}

\cleardoublepage

\tableofcontents

\cleardoublepage
\section{Einführung}

Das folgende Dokument beschreibt das Corporate Design des Vereins Leipzig eSports e.V. und dient als Grundlage für die Gestaltung aller Dokumente die den Verein repräsentieren.
Um den Gestaltenden mehr Freiheit zu lassen, sind die Richtlinien in diesem Dokument so formuliert, dass sie möglichst viel Freiraum für Interpretation lassen.
Lediglich an einigen Stellen sind bindende Regeln definiert und durch die Formulierungen \emph{müssen} oder \emph{nicht dürfen} gekennzeichnet.
Von diesen darf nicht abgewichen werden.

\cleardoublepage
\section{Der Name und das Kürzel}

Der Verein tritt in der Regel mit dem vollen Vereinsnamen \emph{Leipzig eSports e.V.} auf.
Zur Verbesserung des Leseflusses kann das \emph{e.V.} in nicht-formellen Texten wie Artikeln, Blog-Einträgen oder Forenposts weggelassen werden.

Das Kürzel für Clan- oder Gruppenstrukturen in Spielen, Netzwerken, u.ä. lautet immer \emph{LESeV}.
Aus technischen Gründen (Mindestlänge, bereits vergeben) kann es hier notwendig sein, dass von diesem Kürzel abgewichen wird.
Ansonsten sind andere Schreibweisen oder Abkürzungen nicht zulässig.

\cleardoublepage
\section{Das Logo}

Das Logo zeigt den vollen Schriftzug \emph{Leipzig eSports} in gelb und blau und den stilisierten Löwen in orange in einer Wappenform.
\begin{figure}[H]
\includegraphics[width=\textwidth]{Docs/Logo.eps}
\caption{Logo mit Schriftzug}
\end{figure}

Die Schriftart ist Helvetica Neue, bei der das L abgewandelt ist, um der Silhouette des City-Hochhauses zu ähneln.

\begin{figure}[H]
\centering
\includegraphics[width=0.33\textwidth]{Docs/logo_detail_L.png}
\caption{Detailansicht L}
\end{figure}

\subsection{Varianten}

Die erste Variante ist das Logo ohne Schriftzug, nur als Wappen.
Es bietet sich an für Icons und Dokumente in denen das volle Logo aus Platzgründen nicht platziert werden kann.

\begin{figure}[H]
\centering
\includegraphics[width=0.33\textwidth]{Docs/Emblem.eps}
\caption{Logo ohne Schriftzug}
\end{figure}

Als zweite Variante gibt es den Schriftzug ohne Wappen.
Er sollte dann verwendet werden, wenn das Wappen bereits in der Nähe platziert wurde, der Schriftzug aber zusätzlich nötig ist.
Er darf nicht auf eine Art und Weise zusammen mit dem Wappen platziert werden, die den Anschein erweckt, dass diese Anordnung das volle Logo darstellt.

\begin{figure}[H]
\centering
\includegraphics[width=0.66\textwidth]{Docs/Type.eps}
\caption{Logo ohne Wappen}
\end{figure}


\subsection{Farbvarianten} \label{farbvarianten}

Die Farben des Logos dürfen nicht verändert werden. Insbesondere dürfen Farben von Schriftzug und Wappen nicht vertauscht werden (z.B. blaues Wappen, "eSport" in Orange). Es gibt einfarbig schwarze und weiße Varianten des Logos, die mit variabler Deckkraft platziert werden können, jedoch nicht eingefärbt werden dürfen.

% beispielbilder

\subsection{Platzierung}

Um das Logo herum ist ein Abstand einzuhalten der mindestens $\sfrac{1}{7}$ der Höhe des Logos entspricht.\footnote{Beim Schriftzug ohne Wappen ist dieser Abstand entsprechend dem, den das Logo hätte, wenn das volle Logo in gleicher Größte verwendet worden wäre.}
Aufgrund der transparenten Bereiche ist das Platzieren auf Hintergründen mit mehreren Farben oder starken Kontrasten nicht sinnvoll.
Das Logo sollte idealerweise auf sehr hellen oder sehr dunklen, einfarbigen Hintergründen platziert werden, um optimalen Kontrast zu garantieren.
Auch die Platzierung auf Fotos ist möglich, wenn diese wenig Kontrast haben oder vorher mit einem Weichzeichnen-Filter bearbeitet werden.

\begin{figure}[H]
\begin{centering}
\hfill
\fbox{\includegraphics[width=0.3\textwidth]{Docs/logo_placement_bright.png}}
\fbox{\includegraphics[width=0.3\textwidth]{Docs/logo_placement_medium.png}}
\fbox{\includegraphics[width=0.3\textwidth]{Docs/logo_placement_dark.png}}\hfill

\hfill
\fbox{\includegraphics[width=0.3\textwidth]{Docs/logo_placement_picture.png}}
\fbox{\includegraphics[width=0.3\textwidth]{Docs/logo_placement_blur.png}}\hfill
\end{centering}
\caption{Beispiele für Kontrast auf verschiedenen Hintergründen}
\end{figure}

\subsection{Vereinsabkürzung}

Das Logo trägt im Normalfall noch die Abkürzung \emph{e.V.}.
Sie wird hochgestellt hinter dem \emph{eSport} platziert und ist je nach Hintergrund schwarz oder weiß, entsprechend der regulären Schriftfarbe. %schriftart?

%hier Bilder von Logo mit e.V.


\cleardoublepage
\section{Die Farben}

\subsection{Primärfarben}
\begin{samepage}
\cfield[LESOrange]
Orange
\begin{itemize}
\item HTML \#FF7929
\item RGB 255/121/41
\item HSV 22/84/100
\item CMYK 0/60/92/0
\end{itemize}
\end{samepage}

\begin{samepage}
\cfield[LESBlue]
Blau
\begin{itemize}
\item HTML \#2E72C9
\item RGB 46/115/201
\item HSV 214/77/79
\item CMYK 76/41/0/2
\end{itemize}
\end{samepage}

\begin{samepage}
\cfield[LESYellow]
Gelb
\begin{itemize}
\item HTML \#FFD129
\item RGB 255/209/41
\item HSV 47/84/100
\item CMYK 0/18/87/0
\end{itemize}
\end{samepage}

Die Primärfarben orientieren sich an den bisher benutzten Vereinsfarben und sind nur leicht abgeändert, um etwas kräftiger zu wirken.
Sie finden sich nicht nur im Logo sondern auch in hervorhebenden Gestaltungselementen.

\subsection{Sekundärfarben}
\begin{samepage}
\cfield[SECLight]
Helles Grau
\begin{itemize}
\item HTML \#b4b8bd
\item RGB 180/184/189
\item HSV 213/5/74
\item CMYK 3/1/0/29
\end{itemize}
\end{samepage}

\begin{samepage}
\cfield[SECMedium]
Mittleres Grau
\begin{itemize}
\item HTML \#555555
\item RGB 85/85/85
\item HSV 0/0/33
\item CMYK 0/0/0/67
\end{itemize}
\end{samepage}

\begin{samepage}
\cfield[SECDark]
Dunkles Grau
\begin{itemize}
\item HTML \#333333
\item RGB 51/51/51
\item HSV 0/0/20
\item CMYK 0/0/0/80
\end{itemize}
\end{samepage}

Die Sekundärfarben dienen hauptsächlich als Hintergrund- oder Füllfarben.
Das helle Grau hat eine leichte Blaukomponente und ist eine hellere, entsättigte Version des Vereinsblaus.
Das dunkle Grau bildet einen guten Kontrast zu Logo und Gestaltungselementen und sollte für diese als Hintergrund verwendet werden.
Das mittlere Grau ist sowohl als Hintergrund von Gestaltungselementen, als auch als Übergang oder Rand für Elemente in den anderen Grautönen zu benutzen.

\cleardoublepage
\section{Die Schriften}

\subsection{Titillium Web}
\begin{samepage}
Titillium Web ist die Standardschrift, die für alle Texte genutzt wird.

\fonttest{Titillium Web}
\end{samepage}

\subsection{Kanit}
\begin{samepage}
Kanit wird für alle Überschriften und Titel und ausschließlich in fett verwendet.

\fonttest{Kanit}
\end{samepage}

\subsection{Schriftgrößen}

\subsection{Schriftfarben}
Text wird generell in weiß oder schwarz gesetzt und nie in Primärfarben. Generell sollte auf gute Lesbarkeit geachtet werden. So sollte Text auf hellem rau in schwarz, auf mittlerem und dunklem Grau in weiß gesetzt werden.

\cleardoublepage
\section{Gestaltungselemente}

\subsection{Hintergründe}
Hintergründe sollten so wenig Kontrast wie möglich aufweisen, um nicht mit Inhalten und Gestaltungselementen zu kollidieren.
Es bieten sich flache Hintergründe in den Sekundärfarben und weichgezeichnete oder entsättigte Fotos an.
Für besondere Veranstaltungen oder thematisch speziell einzuordnende Dokumente, beispielsweise spielspezifische Streamoverlays, können jedoch auch andere, kontrastarme Hintergründe verwendet werden.

% beispielgrafiken hier

\subsection{Der Winkel}
Alle Elemente deren Linien nicht parallel zu den Dokumenträndern verlaufen, sollten einen Winkel von etwa 35° haben.
Dieser Winkel entspricht dem der Diagonalen durch das Auge des Löwen im Wappen.\footnote{Der Winkel im Auge entspricht nicht exakt 35°, sondern weist einige Dezimalstellen auf. Der Einfachheit halber wird hier jedoch gerundet.}

%illustrationen/grafiken hier

\subsection{Streifen}
Die Streifen dienen als Elemente zur Auflockerung von Hintergründen oder Übergängen.
Sie sind um 35° gedreht und von beliebiger Breite und Anzahl und in beliebiger Farbe.\footnote{Sonstige Regeln bezüglich Kontrast usw. gelten natürlich trotzdem.}

\begin{figure}[H]
\begin{center}
\includegraphics[width=0.5\textwidth]{Docs/streifen_beispiel.png}
\end{center}
\caption{Verschiedene Streifen auf einfachem Hintergrund}
\end{figure}

\begin{figure}[H]
\begin{center}
\includegraphics[width=0.5\textwidth]{Docs/streifen_foto.png}
\end{center}
\caption{Weiße Streifen mit 20\% Deckkraft auf Fotohintergrund}
\end{figure}

\cleardoublepage
\section{Business}

Für geschäftliche Dokumente wie Verträge, Briefe aber auch Visitenkarten oder offizielle Einladungen gelten besondere Regeln.
Hier sollte auf die Verwendung der meisten Gestaltungselemente und Farben verzichtet werden.
Idealerweise verwenden diese Dokumente nur Schwarz, Weiß und Orange mit möglichst viel Weißraum und großen Flächen.

\begin{figure}[H]
\begin{center}
\includegraphics[width=0.5\textwidth]{Docs/Visitenkarte.eps}
\end{center}
\caption{Visitenkarte als Beispiel für businessorientiertes Design}
\end{figure}

\cleardoublepage
\section{Ausnahmen und Abweichungen}

Für bestimmte Medien darf von diesem Corporate Design Manual abgewichen werden.
Diese sind unter anderem:

\begin{itemize}
\item Animationen und Intros
\item Interaktive Designelemente
\item Künstlerische Darstellungen wie z.B. Wallpaper
\item Textilien wie Trikots oder T-Shirts
\end{itemize}

Die Abweichungen müssen sich jedoch auf das jeweilige Werk beschränken und sollten innerhalb des Werkes einheitlich sein.
So ist z.B. ein 3D-Logo für ein Streamintro nur in diesem zu nutzen und nicht auf dem Trikot oder einem Flyer.
Insgesamt sollten Abweichungen dieser Art allerdings auf ein Minimum begrenzt werden. Sie sollten sich ästhetisch an dieses Corporate Design Manual anlehnen.

\subsection{Ausnahmen beim Logo}
Das Logo darf weder in Form noch in Farbe abgeändert werden. Es gelten weiterhin die Regeln aus Abschnitt \ref{farbvarianten}.

\end{document}
